\documentclass[11pt]{article}
% Shared preamble of the RMC kernel write-ups (\input{preamble} after \documentclass).
\usepackage[margin=1in]{geometry}
\usepackage{amsmath,amssymb,bm}
\usepackage{graphicx}
\usepackage{booktabs}
\usepackage{hyperref}

\newcommand{\dd}{\,\mathrm{d}}
\newcommand{\Ein}{E}
\newcommand{\Eout}{E'}
\newcommand{\Ecm}{E_{c}}
\newcommand{\mucm}{\mu_{c}}
\newcommand{\mulab}{\mu_0}

\newcommand{\acespec}[1]{J.~L.~Conlin and P.~Romano,
\href{https://mcnp-green.lanl.gov/pdf_files/TechReport_2019_LANL_LA-UR-19-29016_ConlinRomano.pdf}
{\emph{A Compact ENDF (ACE) Format Specification}}, LA-UR-19-29016 (2019),
Sec.~#1}
\newcommand{\ENDFspec}[1]{\href{https://www.nndc.bnl.gov/endfdocs/ENDF-102-2023.pdf}
{\emph{ENDF-6 Formats Manual}} (2023), #1}

% Common closing section: how the verification figures are produced
\newcommand{\verificationmethod}{%
Each figure is produced by \texttt{verification/verify\_laws.py}. For an incident
energy $\Ein$ along $+z$, $10^6$ collisions are sampled with MC/DC's own reaction
sampler (\texttt{sample\_elastic\_scattering}, \texttt{sample\_inelastic\_scattering}
or \texttt{sample\_fission}), and every emitted neutron is histogrammed in
$48 \times 24$ bins of $(\Eout, \mulab)$, with $\mulab$ the lab cosine to $+z$. The
inverted kernel used by RMC is integrated over the same bins (Gauss--Legendre panels;
two-body lines are split where $\mulab^*(\Eout)$ crosses a cosine edge), multiplied by
the number of collisions, and compared with a Pearson-type $\chi^2$ diagnostic
over bins with at least 5 expected counts (the rest pooled into one bin).
The existing tables list the script's $\chi^2/\mathrm{dof}$ and nominal $p$-value,
the emitted neutrons per collision (sampled vs.\ the integral of the
inverted kernel), and the largest relative change of a bin integral when the
quadrature panels are doubled. Left panels: outgoing-energy marginal; middle:
lab-cosine marginal; right: per-bin $z = (\text{sampled} - \text{inverted}) /
\sqrt{\text{inverted}}$.

\noindent These figures check agreement with the MC/DC sampler, not independent
physical validity. Their existing $p$-values are not calibrated goodness-of-fit
probabilities: the script uses $B$ degrees of freedom, whereas a fixed-total
multinomial histogram with $B$ categories and no fitted parameters uses $B-1$.
Random multiplicities and spectrum assignments can also change the collision-level
covariance from independent Poisson counts; $z$ is then only a count diagnostic.
See \href{https://docs.scipy.org/doc/scipy/reference/generated/scipy.stats.chisquare.html}
{the Pearson test's sampling assumptions}.
Energy bins use the sample extrema, so unsampled tails are excluded and the
reported inverted yield is an integral over that observed window. A formal test
needs independently specified bins, explicit tail categories, and a covariance or
resampling model appropriate to independent collisions. Zero expected density
with observed emissions must be reported rather than discarded. Doubling quadrature
panels gives a convergence indicator, not a rigorous error bound. The stored
figures and tables have not been regenerated for these qualifications.}

\title{Characteristic scattering-source trials for three-dimensional RMC}
\date{}
\begin{document}
\maketitle

\section{Purpose and scope}
The research objective is a general construction with a better total-cost figure
of merit than standard Monte Carlo, not an isolated favorable example.
The prototype in \texttt{examples/rmc/streaming\_source} tests a prespecified
matrix of optical thickness, scattering ratio, source localization and material
contrast. It is steady, one-group, isotropically scattering transport on a
structured Cartesian mesh with vacuum boundaries. Its transport kernel is shared
by both methods. It is not a replacement for MC/DC's continuous-energy solver.

\section{Representation and residual}
Let $q(\mathbf x)$ be total isotropic external emission density,
$T=\boldsymbol\Omega\cdot\nabla+\Sigma_t(\mathbf x)$, and
$S\psi=\Sigma_s\phi/(4\pi)$ with $\phi=\int\psi\,\mathrm d\Omega$.
Represent an approximate collision source by a cellwise constant total emission
density $Q$. Define
\begin{equation}
 \psi_t=T^{-1}\frac{q+Q}{4\pi},\qquad
 r=q/(4\pi)-(T-S)\psi_t
   =\frac{\Sigma_s\phi_t-Q}{4\pi}.
\end{equation}
The trial satisfies streaming, attenuation and vacuum conditions through
characteristics, with no artificial jumps in the trial flux at material faces.
Along a backward ray the scalar-direction integral is
\begin{equation}
 I(\mathbf x,\Omega)=\int_0^{s_b}(q+Q)(\mathbf x-s\Omega)
 \exp\left[-\int_0^s\Sigma_t(\mathbf x-u\Omega)\,\mathrm du\right]\mathrm ds.
\end{equation}
It is integrated exactly for piecewise constant coefficients: a segment of length
$d$ contributes incoming attenuation times
$(q+Q)(1-e^{-\Sigma_t d})/\Sigma_t$, with the continuous limit $(q+Q)d$ in voids.
Then $\phi_t=\mathbb E_{\Omega\sim U(S^2)}I$.

\section{A computable frozen source}
Let $P_{ij}$ map isotropic emission density in cell $j$ to average uncollided
scalar flux in cell $i$. Forward ray integration from uniform positions and
directions in each source cell estimates $P$ without a reference transport solve.
Solve the small source system
\begin{equation}
 (I-\mathop{\rm diag}(\Sigma_s)\widehat P)Q
 =\mathop{\rm diag}(\Sigma_s)\widehat Pq.
\end{equation}
Every sample obeys $\sum_i V_i\Sigma_{t,i}\widehat P_{ij}\le V_j$.
Thus $\mathop{\rm diag}(\Sigma_s)\widehat P$ is a contraction in the volume-weighted
$\ell^1$ norm when $\max\Sigma_s/\Sigma_t<1$ over nonvoid material cells
(voids contribute zero scattering). Positivity and invertibility do not
depend on statistical agreement of $\widehat P$ with $P$. Errors in $Q$ affect
efficiency, not the exact residual identity. Source construction is frozen before
independent final estimates. The prototype stores a dense matrix; sparse or
matrix-free source solves are required before scaling to large meshes.

\section{Randomized baseline and residual estimates}
The cell-average trial tally is not known analytically. Estimate it using
independent scrambled Sobol integrations over position (three coordinates) and
direction (two coordinates), separately in every cell. Each scramble is unbiased;
the variance of their mean is estimated across independent scrambles.
The same local Sobol coordinates may be shared across cells to amortize generator
construction. Cell estimates are then correlated within a scramble; independence
across scrambles is retained. The per-cell variances and mean of per-cell relative
variances remain valid. A combined response must use the covariance or its
scramble-level totals, not a sum of independent-cell variances.

At each residual birth position estimate $\phi_t$ with $K$ spherical directions.
For independent uniform random shifts $U,V$ use
\begin{equation}
 \mu_k=2\{(k+1/2)/K+U\}-1,\qquad
 \varphi_k=2\pi\{k(\sqrt5-1)/2+V\}.
\end{equation}
Each direction is marginally uniform, so $\widehat\phi_t=K^{-1}\sum_k I_k$
is unbiased for every finite $K$. Directions are correlated within the inner
estimate; shifts are independent across residual histories. No deterministic
angular interpolation is substituted for the residual.

Choose a cell with fixed probability $p_i>0$, a position uniformly inside it,
and an independent isotropic outgoing direction. The total source weight is
\begin{equation}
 w=\frac{V_i}{p_i}(\Sigma_{s,i}\widehat\phi_t-Q_i).
\end{equation}
The initial proposal is uniform in volume. It neither accepts nor rejects a birth
using its noisy residual value. Analog transport with signed weights estimates
the cell response to this random source. Conditional expectation over the inner
rays gives the exact residual source. The final estimator is the independent
trial-tally estimate plus the transported residual response. Conditional on $Q$
its expectation is the true scalar flux. Independent source-construction noise
therefore adds no variance through a random conditional mean.

\paragraph{Antipodal inner rays.}
For even $K$, an optional inner estimator draws $K/2$ randomized lattice
directions and evaluates both $\Omega$ and $-\Omega$. Each direction remains
marginally uniform, while the average removes all odd angular components of the
ray integrand exactly. The ray count and expectation are unchanged. This can
reduce inner variance when odd components dominate; it does not guarantee a
lower variance for every angular distribution. Pair covariance is included by
the outer residual-history statistics. It changes neither the residual proposal
support nor the analog transport kernel.

\section{Error and cost}
\paragraph{Uncollided/collided control.}
Setting $Q=0$ gives the simpler hybrid split with trial $T^{-1}q$ and residual
$ST^{-1}q$. The unfitted control uses the same randomized baseline, inner
integration, proposal pilot and transport, but skips the collision-source solve.
It still builds the first-flight response matrix for pilot response scales.
This comparison separates gains due to the fitted RMC source from gains due to
randomized quadrature or analytically integrated uncollided transport alone.
Efficiency against analog SMC does not by itself establish an advantage over
this or other optimized hybrid/variance-reduced baselines.

\paragraph{Forward baseline and frozen response proposal.}
Backward ray estimates can have high variance when a source occupies a small
region. An equivalent forward estimate samples positions and directions in
each emitting cell, integrates attenuated track lengths, and accumulates
$\widehat P(q+Q)$ on all destination cells. Independent scrambles give its
variance. Both baseline estimators are unbiased; compare them on the entire
case matrix, including their cost.

For the correction, a separate pilot estimates the second moment $v_i$ of
the vector of transported cell responses from birth cell $i$, normalized by
approximate cell fluxes from the source solve. Freeze
$p_i=0.9\sqrt{v_i}/\sum_j\sqrt{v_j}+0.1V_i/\sum_jV_j$ for final histories.
Pilot draws are discarded from final estimates, with independent final seeds.
The uniform mixture preserves support even when the pilot misses a contribution.
The response normalization needs no reference SMC solve. This is an efficiency
heuristic for the aggregate relative-variance objective, not an optimality claim
for every individual response or for unequal history costs. Pilot cost is part
of source/setup time.

\paragraph{Source-directed angular importance.}
Another unbiased baseline samples an emitting cell $j$ with probability
$a_j\propto |q_j+Q_j|V_j$. At a uniform target point, its bounding sphere
subtends a cone containing every ray that can intersect the emitter. Sample
uniformly in that cone (the full sphere when the point lies inside the bounding
sphere), integrate only emitter $j$ along the attenuated ray, and multiply by
$\Delta\Omega_j/(4\pi a_j)$. Summing expectations over emitters gives the same
trial scalar flux. This avoids requiring an unconstrained random direction to
find a small remote source. Ray attenuation still includes all traversed
materials; no attenuation or emission is replaced by a cell-average optical
length. Six-dimensional scrambled Sobol points provide position, emitting-cell
selection and cone direction. The finite cone contains the complete emitting
cell, so no physical contribution is excluded.

\paragraph{Separating external and collision emission.}
Linearity permits independent estimates of $T^{-1}q$ and $T^{-1}Q$.
The split experiment uses source-directed integration only for localized external
emission, and ordinary direction integration for the smoother collision source.
Its fixed geometry rule is to split when fewer than half the cells emit external
particles; otherwise use the original combined backward estimator. This is a
heuristic, not a general optimality result. Shared randomness within a replicate
is allowed, but the variance is computed from the combined replicate values,
including their covariance. All extra ray and quadrature construction costs are
included. The rule uses no measured SMC results or case-specific FOM selection.

Add the variance of the independently scrambled baseline mean to the
residual-history variance. The latter includes the inner-ray noise. In particular,
\begin{equation}
 \mathbb E[\widehat r^2\mid\mathbf x]
 =r(\mathbf x)^2+\operatorname{Var}(\widehat r\mid\mathbf x).
\end{equation}
Increasing $K$ reduces only the second term and may worsen total FOM. Independent
complete-run replicates check the reported uncertainty and compare empirical
variance times total wall time. Times include source construction, linear solve,
baseline quadrature, residual evaluation and transport, each once. JIT compilation
is warmed consistently and its separate cold cost is recorded. FOMs use the same
cell responses and normalization for both methods; all matrix cases, including
losses, are retained.

\section{Generalization and verification gates}
\paragraph{Linear collision-source reconstruction.}
The linear-source experiment retains the flat-source system for cell means,
then reconstructs slopes from finite differences of its approximate scalar flux:
$Q_i(\mathbf t)=Q_{i0}+\sum_a Q_{ia}t_a$, with
$Q_{ia}=\Sigma_{s,i}(h_a/2)\partial_a\phi$ and $t_a\in[-1,1]$.
Scale all slopes in a cell by
$\min(1,0.95Q_{i0}/\sum_a|Q_{ia}|)$ to retain nonnegative emission.
Reconstructing scalar flux before multiplying by each cell's cross section
avoids differentiating an artificial collision-source material jump.
This is a frozen heuristic reconstruction, not a higher-order Galerkin solve.
It is nevertheless a valid trial: integrate its affine source along each ray
segment analytically using $I_0=\int_0^d e^{-\Sigma_t s}\mathrm ds$ and
$I_1=\int_0^d s e^{-\Sigma_t s}\mathrm ds=(I_0-de^{-\Sigma_t d})/\Sigma_t$.
A small-optical-depth series avoids cancellation in $I_1$.
Evaluate the same reconstructed $Q(\mathbf x)$ in the residual weight. Thus its
source reconstruction error affects variance, not expectation. The current
linear experiment uses only the backward baseline; other baseline combinations
are explicitly rejected until their polynomial-source integration is implemented.

The construction does not require homogeneous cells beyond the ray integrator's
ability to integrate source and attenuation. Anisotropic scattering requires
direction-dependent $Q$ and the corresponding incident angular integral.
Continuous energy requires an incident-energy integral, energy-dependent ray
attenuation, and multiplicities consistent with the production transport kernel.
Reflective boundaries require continued reflected characteristics with controlled
termination; multiplying media need a separate stability argument. None of these
extensions is verified by the present one-group experiments.

Component checks cover exact heterogeneous-ray attenuation, first-flight
conservation, residual support and pure absorption. Integrated checks compare
scattering runs with independent SMC and verify error coverage over full-run
replicates. Mesh growth, optical thickness and source contrast are separate
performance axes; selecting a single favorable configuration is not acceptance.
\end{document}
